\section{Datasets and Training Setup}
\label{app:datasets}

\begin{table}[t]
\centering
\small
\begin{tabular}{lrr}
\toprule
 & PTB & WikiText-2 \\
\midrule
Vocabulary (types)          & 10{,}000  & 33{,}278 \\
Train tokens                & 929{,}589 & 2{,}088{,}628 \\
Validation tokens           & 73{,}760  & 217{,}646 \\
Test tokens                 & 82{,}430  & 245{,}569 \\
\midrule
Validation blocks           & 1{,}152   & 3{,}400 \\
Scored validation tokens    & 72{,}576  & 214{,}200 \\
Test blocks                 & 1{,}287   & 3{,}837 \\
Scored test tokens          & 81{,}081  & 241{,}731 \\
\midrule
Unigram val.\ ppl.          & 688.82    & 964.82 \\
\quad slot $0$ included     & 687.45    & 965.37 \\
Bigram val.\ ppl.           & 445.24    & 692.50 \\
\bottomrule
\end{tabular}
\caption{The two corpora as the loader sees them. Types include \texttt{<eos>} and are counted
on the training split; no validation or test type is missing from it. Blocks are the
non-overlapping $\seqlen=64$-token evaluation blocks; $63$ tokens per block are scored (slot $0$
is dropped for every model). The bigram is add-one smoothed towards the unigram.}
\label{tab:datasets}
\end{table}

Both corpora are word-level and pass through one reader, one tokenisation function and one
evaluation protocol. All tuning is on Penn Treebank; WikiText-2 serves only for the transfer
test of \Cref{sec:exp-wt2}. \Cref{tab:datasets} lists counts measured on the token streams the
runs consume.

\paragraph{Penn Treebank.}
We use the standard word-level language-modelling version of the Penn Treebank
\citep{marcus1993building}, with the preprocessing of \citet{mikolov2012statistical} as
distributed: lower-cased, numbers replaced by \texttt{N}, rare
words replaced by \texttt{<unk>}, punctuation removed, one sentence per line. The only type with
an upper-case
letter is \texttt{N}, which accounts for 32{,}481 of the 929{,}589 training tokens. Only 46
training types contain a digit, and all are compounds, decades or ordinals (\textit{10-year},
\textit{1920s}, \textit{19th}). The only types with neither a letter nor a digit are
\texttt{\#}, \texttt{\$}, \texttt{\&}, the apostrophe and two escaped asterisks. The reader
appends \texttt{<eos>} to every line, so the sentence boundary is a token the model must
predict. There are 42{,}068 training sentences, 3{,}370 validation and 3{,}761 test sentences,
and \texttt{<eos>} is $4.57\%$ of the validation tokens. The vocabulary is built from the
training split alone: $9{,}999$ word types plus \texttt{<eos>}, the standard $10{,}000$. No
validation or test type is missing from it, so the reader's \texttt{<unk>} fallback never fires
and the unigram baseline needs no smoothing.

\paragraph{WikiText-2.}
We use the word-level release of \citet{merity2017pointer} (\texttt{wikitext-2-v1}, not the
\texttt{-raw-} release, whose preprocessing differs). Its files are already mapped to a closed
vocabulary, and \texttt{<unk>} accounts for $54{,}625$ training tokens. They are read by the same
function as PTB, with the same whitespace split, one \texttt{<eos>} per line and a vocabulary
built from the training split: $33{,}278$ types including \texttt{<eos>}, again with no
validation or test type missing. The files keep the article structure. Of the 3{,}760
validation lines, 1{,}299 are blank and 620 are headings (\texttt{= Heading =},
\texttt{= = Sub = =}, and so on, up to five levels). We score blank lines and headings as they
are, so that the token count matches the released files; a blank line contributes a bare
\texttt{<eos>}. On the validation split blank lines are $0.60\%$ of the tokens and heading lines
$2.37\%$. Filtering the headings would change the denominator of the perplexity and make our
numbers incomparable with published ones.

\paragraph{Batching and slot alignment.}
Each split is one \texttt{<eos>}-joined stream, cut into blocks of $\seqlen = 64$ tokens as in
nanoGPT \citep{karpathy2022nanogpt}, and every model sees the same blocks in the same order.
Training draws blocks at uniformly random offsets, so a block usually starts mid-sentence. The
causal PT's $\text{logits}[:,t]$ predicts $w_t$ from $w_{\prefix{t}}$ (\Cref{sec:output}), so
it scores all $\seqlen$ tokens of a block, including $w_0$ from \textsc{root} alone. A causal
transformer predicts only $w_1,\dots,w_{\seqlen-1}$. To score the same token set, the
transformer wrapper runs the backbone on $w_{0:\seqlen-2}$ and shifts its logits right by one
slot, filling slot $0$ with zeros. Both models then drop slot $0$ in training and in evaluation,
so neither receives gradient from a slot the other cannot predict, and the wrapper's loss
asserts that the zero-filled slot is never scored. Both models are thus asked for
$p(w_t \mid w_{\prefix{t}})$ for $t = 1,\dots,\seqlen-1$: $63$ scored tokens per block,
$72{,}576$ on PTB validation and $214{,}200$ on WikiText-2 validation.

\paragraph{Evaluation.}
The evaluation loop walks fixed, non-overlapping blocks in corpus order over the whole split, in
batches of $16$, and samples nothing. Trailing tokens that do not fill a block are dropped: $53$, $32$
and $62$ tokens on the three PTB splits and $52$, $46$ and $1$ on WikiText-2, at most $0.08\%$
of a split. Perplexity is the exponential of the mean negative log-likelihood over all scored tokens of the
split, not a mean of per-block perplexities. Training perplexity is measured the same way on a fixed $320$-block slice of the
training split.

\paragraph{Unigram baselines.}
The unigram reference of \Cref{sec:tasks} is the maximum-likelihood unigram of the training
split. It is scored through the same loop, so that it has the same $63$ tokens per block as
every model: $688.82$ on PTB validation and $964.82$ on WikiText-2 validation. Including slot $0$
gives $687.45$ and $965.37$, a difference of at most $0.2\%$, in opposite directions on the two
corpora. A model that does not clear the unigram by a wide margin has not learned to use
context, whatever its loss curve looks like.

\paragraph{WikiText-2 is not a second tuning corpus.}
Every hyperparameter (learning rate and schedule, width, iteration count, temperature and
damping) is selected on PTB and reused unchanged on WikiText-2; no grid is rerun there. This is
what makes \Cref{sec:exp-wt2} a transfer test. \citet{kuang2026scaling} make no claim that their
transfer method holds across corpora, so this test is not covered by their results. Two things change with the corpus: the vocabulary
grows from $10{,}000$ to $33{,}278$ types, and the training set from $0.93$M to $2.09$M tokens.
The word--label factor is $|\Vocab| \times \nlab$, so the parameter count changes even at fixed
$\nlab$, and WikiText-2 results are compared only with baselines trained on WikiText-2 through
the same loop. The readout bound of \Cref{prop:readout} also binds harder, with $3.3$ times the
vocabulary to place in an affine family of dimension $\nlab$. It binds on a transformer's output
head in the same way at matched width (\Cref{sec:cmp-readout}), but a weaker absolute number on
WikiText-2 than on PTB partly reflects it.

\paragraph{Training protocol.}
Every model is trained by one loop, which reaches model-specific behaviour only through three
methods: the loss on a scored token set, an $L_2$ regulariser on the arc scores (zero for a
transformer) and a parameter count. PTB has under a million training tokens, so $15{,}000$ steps
of $16$ blocks is $16.5$ epochs, and the larger baselines reach their validation optimum long
before the end (a 1.2M-parameter transformer at step $5{,}000$, against $15{,}000$ for the
causal PT). Every model is therefore scored at its own best-validation checkpoint, and test
perplexity is measured on that checkpoint. Dropout is $0$ for every model: \citet{wu2023probabilistic} do not
specify where it would enter the update equations, and any placement would add a map with no
factor behind it (\Cref{sec:variants}). The causal PT's word unary $\bfac$ is set to the log
unigram frequency of the training split and not trained. Runs used NVIDIA A40, TITAN~RTX and
L20 GPUs.

\paragraph{Model configurations.}
The baselines are nanoGPT decoders with four layers and four heads at widths $\dmodel$ from
$16$ to $256$; the looped transformer applies one such block four times with shared weights.
Unless stated otherwise the causal PT uses $\nchan=2$ channels, $\niter=3$ iterations, clip
$\clip=3$, $\rho=2$ predictive rounds (\Cref{app:full-updates}), the mean-field readout,
initialisation scale $0.02$ and a learning rate of $2\times10^{-2}$. Following the step size of
Wu and Tu's Appendix~B.1, the label update can be damped,
$\QZ{i}^{(t)} = \alpha_Z Q^{\star}_{Z_i} + (1-\alpha_Z)\QZ{i}^{(t-1)}$, where $Q^{\star}_{Z_i}$
is the undamped update and $t$ indexes iterations; $\alpha_Z=1$ means no damping. This differs
from their message damping (their Appendix~B.2). The scaling ladders, the gain and rank grids
and the global-variable grid run undamped under the standardised temperature; the constants,
depth and readout sweeps use the damped configuration ($\alpha_Z=0.25$, original temperature,
$\nlab=\krank=32$); the factored-label grid uses $\alpha_Z=0.25$ and $\nchan=4$ at the original
temperature; the switch ladder runs at $\dmodel=\nlab=64$.
